\begin{otherlanguage}{english}
\chapter*{Abstract}
\addcontentsline{toc}{chapter}{Abstract}

Test-time compute can improve mathematical reasoning by extending one
generation trajectory or sampling multiple independent trajectories, but a
fixed configuration may waste computation on easy inputs and remain
insufficient for difficult ones. This thesis studies adaptive allocation for a
small open-weight reasoning model without an external verifier and accounts
for cost using realized generated tokens.

Using Qwen3-0.6B, we freeze dataset-specific action spaces, policies, and
disjoint 100-problem Test partitions for GSM8K and MATH-500. Fixed,
breadth-adaptive, token-cap-adaptive, and joint-heuristic policies are replayed
against a common generation pool. Primary outcomes are exact-answer accuracy
and realized token cost, supported by Wilson intervals, paired McNemar tests,
latency, validity, and truncation analysis.

On GSM8K, agreement stopping at $L=1024$ reaches 76\% accuracy with 1,994.52
mean generated tokens, matching fixed $k=4,L=1024$ while using 21.3\% fewer
tokens. On MATH-500, the Validation-selected fixed $k=2,L=4096$ policy reaches
66\%. The joint heuristic reaches only 36\% with 5,788.81 tokens and is
significantly worse than the selected policy. An agreement-stopping comparator
at $L=4096$ reaches 73\%, matching the aggregate accuracy of fixed
$k=4,L=4096$ while using approximately 20\% fewer tokens.

The evidence does not support a single joint allocation rule across both
datasets. Fixed policies selected for each dataset and policies adapting one
axis provide stronger operating points. The primary evidence uses one
generation seed, one model family, and two Test partitions of 100 problems.
Non-thinking mode, the native Qwen3 budget interface, and budget forcing remain
proposed robustness evaluations rather than completed evidence.

\noindent\textbf{Keywords:} adaptive test-time compute; small reasoning models;
breadth--token-cap allocation; efficient reasoning; Qwen3.
\end{otherlanguage}
